\subsection{Countable compactness for regular logics}
\label{sec:p03-general-pinning}

The connection between compactness and well-orderings does not depend on
an embedding into an infinitary language. Relativization and simultaneous
substitution suffice. The latter operation also separates the
interpretations of constants, functions and nullary relations in the
models used in the construction.

\begin{hbconvention}{def:p03-general-pinning-setting}{Scope and cardinality}
In this subsection, \(\mathscr L\) is a one-sorted regular logic in the
sense of \cref{def:p03-w03-p03-047-closure,def:p03-w03-p03-047-relativization}.
Vocabularies may have any set cardinality and may contain constants,
finitary functions and relations, including nullary relations.
Domains are nonempty. Work in classical set theory with Choice.
Use countable compactness and pinning as in
\cref{def:p03-w03-p03-041-compactness,def:p03-w03-p03-041-pinning},
now for these vocabularies and this logic. In particular,
\(w_{\aleph_0}(\mathscr L)\) is the supremum of ordinals pinned by
\emph{at most countably many sentences}; there is no bound on the
size of the vocabulary of those sentences. Pinning requires all models
to give well-orders of the designated relation's field and one model
to give the specified order type, not that every model have that type.
\end{hbconvention}

\begin{hbtheorem}{thm:p03-general-fibres}{proved=annex:p03.general-fibres}{Uniform model fibres}
For any vocabulary \(\tau\), there is an auxiliary vocabulary with a
binary function \(p\) such that, for every \(\varphi\in\mathscr L[\tau]\),
there is a sentence \(\widehat\varphi(d)\) with a distinguished fresh
constant \(d\). Each auxiliary structure \(\mathfrak B\) and element
\(b\in B\) determine a nonempty \(\tau\)-structure \(\mathfrak B_b\)
on
\[
 D_b=\{p^{\mathfrak B}(b,z):z\in B\}
\]
for which
\[
 (\mathfrak B,b)\models_{\mathscr L}\widehat\varphi(d)
 \quad\Longleftrightarrow\quad
 \mathfrak B_b\models_{\mathscr L}\varphi.
\]
The auxiliary vocabulary and the induced structures are the same for
all \(\varphi\). No auxiliary axioms, one for each symbol of \(\tau\),
are required. Any set-indexed family of nonempty \(\tau\)-structures can
be realized, up to isomorphism, as fibres over distinct index elements
in a single auxiliary structure. Its ambient domain may also contain
any prescribed additional set of points, without changing these fibres.
\end{hbtheorem}

\begin{hbtheorem}{thm:p03-general-pinning}{proved=annex:p03.general-pinning}{Countable compactness and well-ordering numbers}
For a one-sorted regular logic \(\mathscr L\), the following are equivalent:
\begin{enumerate}
 \item \(\mathscr L\) is countably compact.
 \item No countable theory of \(\mathscr L\) pins down an infinite ordinal.
 \item \(w_{\aleph_0}(\mathscr L)=\omega\).
\end{enumerate}
If countable compactness fails, countable theories pin down both
\(\omega\) and \(\omega+1\).
This is the regular-logic form of
\cite[Ch.~II, Proposition~5.2.4, p.~62]{Ebbinghaus1985ExtendedLogics}.
\end{hbtheorem}

The general theorem requires no finite-occurrence property, no countable
vocabulary and no prescribed formula grammar. Its proof uses the semantic
closure operations of regularity. The separate
\cref{thm:p03-w03-p03-041-pinning-compactness} remains useful for the
weaker, explicitly syntactic package (F1)--(F4): a sublogic satisfying
that package need not have the full simultaneous-substitution property.
